
\subsection{The integrability of $k_\chi ^T(x,f)$}

For $K$-finite compactly supported $f$ as above, let 
\begin{equation}\label{eq:5.12}
K_{P,\chi }(x,y)=\sum_{P_1\subset P}n_P(A_1)^{-1}
\int_{\Pi _u(\overline{M(\aaa )})}
\sum_{\phi \in {\mathcal{B}}_{P_1}(\pi )_\chi }
E_P(x,I_{P_1}(\pi ,f)\phi )\overline{E_P(y,\phi )}d\pi
\end{equation}
and 
\begin{equation}\label{eq:5.13}
k_\chi ^T(x,f)=\sum_P(-1)^{\dim (A_P/Z)}
\sum_{\delta \in P(\qqq )\backslash G(\qqq )}
K_{P,\chi }(x,x)\tau _P^{\wedge }(H(\delta x)-T).
\end{equation}

\begin{lemma}\label{lem:5.9}
$K_{P,\chi }(x,y)$ is the kernel function for the regular 
representation $R_P(f)$ of $\overline{G(\aaa )}$ on 
\[
L^2(G(\qqq )\backslash \overline{G(\aaa )})_{\mathcal{P},\chi }
=L^2(G(\qqq )\backslash \overline{G(\aaa )})_{\mathcal{P}}
\cap L^2(G(\qqq )\backslash \overline{G(\aaa )})_{\chi},
\]
where $\mathcal{P}$ denotes the associate class of $P$. In particular,
$K_{G,\chi }(x,y)=K_\chi (x,y)$ is the kernel function for the regular
representation $R_P(f)$ of $\overline{G(\aaa )}$ on
$L^2(G(\qqq )\backslash \overline{G(\aaa )})_\chi $.
\end{lemma}

\begin{remark}\label{rem:5.10} In \cite{A15}, different notation is used: $R_P(f)$ denotes
the convolution operator associated to right translation on
\[
L^2(N_P(\aaa )M_P(\qqq )\backslash \overline{G(\aaa )}))
=\oplus_\chi
L^2(N_P(\aaa )M_P(\qqq )\backslash \overline{G(\aaa )}))_\chi .
\]
\end{remark}

\begin{theorem}\label{thm:5.11}
$\sum_{\chi \in \mathcal{X}}K_{P,\chi }(x,y)
=\sum_{{\mathcal O}\in \mathcal{C}}
K_{P,{\mathcal O}}(x,y)$.
\end{theorem}

proof: Both sides represent the kernel function for the regular
representation $R_P(f)$ of $\overline{G(\aaa )}$
on $L^2(G(\qqq )\backslash \overline{G(\aaa )})$. $\Box $

The trace formula is obtained by integrating both sides of 
\[
\sum_{\chi \in \mathcal{X}}\Lambda ^TK_{P,\chi }(x,x)
=\sum_{{\mathcal O}\in \mathcal{C}}
\Lambda ^TK_{P,{\mathcal O}}(x,x) 
\]
over $G(\qqq )\backslash \overline{G(\aaa )}$,
where $\Lambda ^T$ is the
truncation operator defined below. We know from earlier sections that the
right-hand side is a polynomial in $T$, so the left-hand side is too. The 
\textbf{global trace formula} is the constant term of this expression. 

To prove the integrability of $k_\chi ^T(x,f)$ a truncation operator is
introduced.
If $\phi \in C(G(\qqq )\backslash \overline{G(\aaa )})_{gen}$
then define,
as in \cite[p 89]{A3}, the \textbf{truncation operator} by 
\begin{equation}\label{eq:5.14}
\Lambda ^T\phi (x)
=\sum_P(-1)^{\dim (A/Z)}
\sum_{\delta \in P(\qqq )\backslash G(\qqq )}
\int_{N(\qqq )\backslash N(\aaa )}\phi (n\delta
x)dn\cdot \tau _P^{\wedge }(H(\delta x)-T).
\end{equation}
More generally, if $P_1\subset G$ is a parabolic and if
$\phi \in C(P_1(\qqq )\backslash \overline{G(\aaa )})_{gen}$ then
define, as in \cite[p 97]{A3}, the \textbf{truncation operator} by 
\begin{equation}\label{eq:5.15}
\Lambda ^{T,P_1}\phi (x)
=\sum_{R , P_0\subset R\subset P_1}
(-1)^{\dim (A_R/A_1)}
\sum_{\delta \in R(\qqq )\backslash P_1(\qqq )}
\int_{N_R(\qqq )\backslash N_R(\aaa )}
\phi (n\delta x)dn\cdot \tau_R^{\wedge ,1}(H(\delta x)-T).
\end{equation}
In the classical case of $SL(2)$, the truncation operator acting
on the Eisenstein series has the effect of subtracting
off the ``$y^s$ term'' in its Fourier expansion. This allows one 
to compute the inner product of two truncated Eisenstein series
(see Theorem~\ref{thm:5.27} below).

\begin{example}\label{ex:5.12}
For $G=SL(2),P_1=G$ this formula becomes 
\begin{equation}\label{eq:5.16}
\Lambda ^T\phi (x)=\phi (x)-\sum_{\delta \in B(\qqq )\backslash G(\qqq %
)}\int_{N(\qqq )\backslash N(\aaa )}\phi (n\delta x)dn\cdot \tau
_B^{\wedge }(H(\delta x)-T).
\end{equation}
\end{example}

\begin{lemma}\label{lem:5.13}
The truncation operator is a self-adjoint idempotent on
$C(G(\qqq )\backslash \overline{G(\aaa )})_{gen}
\cap L^2(G(\qqq )\backslash \overline{G(\aaa )})$ .
\end{lemma}
proof: This is an immediate consequence of the non-metaplectic
version \cite{A3} since the metaplectic cover splits over $N$. $\Box$

\begin{lemma}\label{lem:5.14}
(\cite[Lemma 1.5]{A3}) For
$\phi \in C(P(\qqq )\backslash \overline{G(\aaa )})_{gen}$, we
have 
\[
\sum_{ P_1 , P_0\subset P_1\subset P }
\sum_{\delta \in P_1(\qqq )\backslash P(\qqq )}
\Lambda ^{T,P_1}\phi (\delta x)\tau _1^P(H_0(\delta x)-T)
=\int_{N(\qqq )\backslash N(\aaa )}\phi (nx)dn.
\]
\end{lemma}

proof in the case $G=SL(2)$: Suppose $P=G$. First, from the fact that 
\begin{equation}\label{eq:5.17}
\sum_{ P_1 , R\subset P_1\subset Q}
 (-1)^{\dim (A_R/A_1)}\tau_R^{\wedge ,1}(X)\tau _1^Q(X)
=\left\{ 
\begin{array}{c}
1,\;\;\;\;Q=R, \\ 
0,\;\;\;\;Q\neq R,
\end{array}
\right.
\end{equation}
(this is Langlands' combinatorial lemma; see \cite{A3})
% Claude (2026): the original cited a nonexistent "lemma 6.9".
, we have $\tau
_B^{\wedge ,B}(X)\tau _B^G(X)=\tau _B^{\wedge ,G}(X)\tau _G^G(X)=\tau
_B^{\wedge }(X)$. The LHS of the equation in the above lemma is 
\begin{eqnarray*}
&&\sum_{ P_1, B\subset P_1\subset G } 
\sum_{\delta \in P_1(\qqq )\backslash G(\qqq )}
\Lambda ^{T,P_1}\phi (\delta x)\tau _1^G(H_0(\delta x)-T) \\
&=&\Lambda ^{T,G}\phi (x)\tau _G^G(H_0(x)-T)+ \\
&&+\sum_{\delta \in B(\qqq )\backslash G(\qqq )}\Lambda ^{T,B}\phi
(\delta x)\tau _B^G(H_0(\delta x)-T) \\
&=&\Lambda ^{T,G}\phi (x)+\sum_{\delta \in B(\qqq )\backslash G(\qqq %
)}\int_{N(\qqq )\backslash N(\aaa )}\phi (n\delta x)dn\cdot \tau
_B^{\wedge ,B}(H_0(\delta x)-T)\cdot \tau _B^G(H_0(\delta x)-T) \\
&=&\Lambda ^{T,G}\phi (x)+\sum_{\delta \in B(\qqq )\backslash G(\qqq %
)}\int_{N(\qqq )\backslash N(\aaa )}\phi (n\delta x)dn\cdot \tau
_B^{\wedge }(H_0(\delta x)-T) \\
&=&\phi (x),
\end{eqnarray*}
by the previous example.

Suppose $P=B$. The LHS of the equation in the above lemma is 
\[
\Lambda ^{T,B}\phi (x)\tau _B^B(H_0(x)-T)
=\int_{N(\qqq )\backslash N(\mathbb{A})}\phi (nx)dn, 
\]
since $\tau _B^B(X)=1$.

This verifies the lemma in both cases. $\Box $

\begin{lemma}\label{lem:5.15}
(\cite[Lemma 2.2]{A3})

\begin{eqnarray*}
k_\chi ^T(x,f)=\sum_{ P_1,P_2 {\rm s.t.} P_0\subset P_1\subset P_2}
\sum_{\delta \in P_1(\qqq )\backslash G(\qqq )}
\sigma _1^2(H_0(\delta x)-T)\times  \\
\times \left\{ \sum_{ P, P_1\subset P\subset P_2 }
(-1)^{\dim (A/Z)}\Lambda _2^{T,P}K_{P,\chi }(\delta x,\delta x)\right\} .
\end{eqnarray*}
(The subscript $2$ in $\Lambda _2^{T,P}$ means that it is to be applied to
the second variable.)
\end{lemma}

\begin{example}\label{ex:5.16}
Let $G=SL(2)$. The above lemma becomes 
\begin{eqnarray*}
k_\chi ^T(x,f)=
\sum_{ P_1,P_2 {\rm s.t.} (P_1,P_2)\in  \{(B,B),(B,G),(G,G)\}} 
\sum_{\delta \in P_1(\qqq )\backslash G(\qqq )}
\sigma_1^2(H_0(\delta x)-T)\times  \\
\ \times \left\{ \sum_{ P , P_1\subset P\subset P_2 }
(-1)^{\dim (A/Z)}\Lambda _2^{T,P_1}K_{P,\chi }(\delta x,\delta x)\right\}  \\
=-\sum_{\delta \in B(\qqq )\backslash G(\qqq )}\sigma _B^B(H_0(\delta
x)-T)\Lambda _2^{T,B}K_{B,\chi }(\delta x,\delta x)\;\;(\mathrm{the}\;(B,B)\;%
\mathrm{term}) \\
+\sigma _G^G(H_0(x)-T)\Lambda _2^{T,G}K_{G,\chi }(x,x)\;\;(\mathrm{the}%
\;(G,G)\;\mathrm{term}) \\
+\sum_{\delta \in B(\qqq )\backslash G(\qqq )}\sigma _B^G(H_0(\delta
x)-T)\times  \\
\;\;\times \left\{ -\Lambda _2^{T,B}K_{B,\chi }(\delta x,\delta x)+\Lambda
_2^{T,B}K_{G,\chi }(\delta x,\delta x)\right\} \;\;\;\;\;\;\;(\mathrm{the}%
\;(B,G)\;\mathrm{term}) \\
=\Lambda _2^{T,G}K_{G,\chi }(x,x)+ \\
+\sum_{\delta \in B(\qqq )\backslash G(\qqq )}\tau _B(H_0(\delta
x)-T)\left\{ -\Lambda _2^{T,B}K_{B,\chi }(\delta x,\delta x)+\Lambda
_2^{T,B}K_{G,\chi }(\delta x,\delta x)\right\}  \\
=\Lambda _2^{T,G}K_{G,\chi }(x,x)+\sum_{\delta \in B(\qqq )\backslash G(%
\qqq )}\tau _B(H_0(\delta x)-T)\Lambda _2^{T,B}K_{G,\chi }(\delta x,\delta
x).
\end{eqnarray*}
by \eqref{eq:2.5}.
% Claude (2026): the original cited "lemma 4.7" here.
The definition is 
\begin{equation}\label{eq:5.18}
k_\chi ^T(x,f)=K_{G,\chi }(x,x)-\sum_{\delta \in B(\qqq )\backslash G(\mathbb{%
Q})}K_{B,\chi }(\delta x,\delta x)\tau _B^{\wedge }(H_0(\delta x)-T).
\end{equation}

We claim now that 
\begin{equation}\label{eq:5.19}
\Lambda _2^{T,G}K_{G,\chi }(x,x)=K_{G,\chi }(x,x)-\sum_{\delta \in B(\qqq %
)\backslash G(\qqq )}\tau _B(H(\delta x)-T)\Lambda _2^{T,B}K_{G,\chi
}(\delta x,\delta x).
\end{equation}
To see this, note 
\[
\Lambda _2^{T,G}K_{G,\chi }(x,x)=K_{G,\chi }(x,x)-\sum_{\delta \in B(\qqq %
)\backslash G(\qqq )}\int_{N(\qqq )\backslash N(\aaa )}\Lambda
_2^{T,B}K_{G,\chi }(x,n\delta x)dn\cdot \tau _B^{\wedge }(H(\delta x)-T).
\]
On the other hand, 
\[
\Lambda _2^{T,B}K_{G,\chi }(x,x)=\int_{N(\qqq )\backslash N(\aaa %
)}K_{G,\chi }(x,nx)dn\cdot \tau _B^{\wedge ,B}(H(\delta x)-T),
\]
so 
\begin{eqnarray}
\sum_{\delta \in B(\qqq )\backslash G(\qqq )}\int_{N(\qqq )\backslash
N(\aaa )}\Lambda _2^{T,B}K_{G,\chi }(x,n\delta x)dn\cdot \tau _B^{\wedge
}(H(\delta x)-T)  \nonumber \\
=\sum_{\delta \in B(\qqq )\backslash G(\qqq )}\Lambda _2^{T,B}K_{G,\chi
}(x,\delta x)\cdot \tau _B^{\wedge }(H(\delta x)-T)  \nonumber \\
=\sum_{\delta \in B(\qqq )\backslash G(\qqq )}\Lambda _2^{T,B}K_{G,\chi
}(\delta x,\delta x)\cdot \tau _B(H(\delta x)-T),
\label{eq:5.20}\end{eqnarray}
since the volume of $N(\qqq )\backslash N(\aaa )$ and $\tau _B(X)=\tau
_B^{\wedge }(X)$ in this case. The claimed result follows.
\end{example}

\begin{example}\label{ex:5.17}
Let $G=SL(2),\;P=G$. Then 
\[
E_G(x,\phi )=E(x,\phi )=\sum_{\delta \in B(\qqq )\backslash G(\qqq %
)}\phi (\delta x)\delta _B(\delta x)^{1/2}.
\]

If $P=B$ then $E_B(x,\phi )=\phi (x)\delta _B(x)^{1/2}$.
\end{example}

\begin{example}\label{ex:5.18}
Let $G=SL(2),\;P=B$. For $\pi \in \Pi _u(\overline{A(\aaa )})$,
$\chi \in \mathcal{X}$,
the set ${\mathcal{B}}_B(\pi )_\chi $ is an orthonormal basis for
${\mathcal{H}}_B(\pi )_\chi $. We have 
\[
{\mathcal{H}}_B(\pi )_\chi =\left\{
\begin{array}{c}
\phi:N(\aaa )A(\qqq )\setminus \overline{G(\aaa )}\rightarrow
{\mathbb C}\ {\rm genuine\ }|\ \\
(1)\ \  \phi {\rm\ is\  right}\  \overline{K_0}\ {\rm -finite},\\
(2) \ \ \phi(ag)=\pi(a)\phi(g),
\ \ a\in \overline{A(\aaa )},\ g\in \overline{G(\aaa )},\\
(3)\ \  ||\phi ||<\infty
\end{array}
\right\}.
\]
\end{example}

\begin{example}\label{ex:5.19}
Let $G=SL(2),\;P=G$. For $\pi \in \Pi _u(\overline{G(\aaa )})$
cuspidal,$\;\chi \in 
\mathcal{X}$, the set ${\mathcal{B}}_G(\pi )_\chi $ is an orthonormal basis
for ${\mathcal{H}}_G(\pi )_\chi $. If $\chi = (G,\pi )$
(in $\mathcal{X}$) then this space is non-zero and consists of
(measurable equivalence classes of) square-integrable
genuine functions on $G(\qqq )\backslash \overline{G(\aaa )}$
which are cuspidal automorphic forms of type $\pi_\infty $.
Since almost every $\pi _p$ is unramified, the action of almost
all $\overline{K_{0,p}}$ is trivial on the
functions in this space. Since the functions in
this space are all $\overline{K_0}$-finite,
we may decompose it by the action of $\overline{K_0}$ on
the right: 
\[
{\mathcal{H}}_G(\pi )_\chi =
\bigoplus_{\sigma \in \Pi(\overline{K_0})}
{\mathcal{H}}_G(\pi )_{\chi ,\sigma }
\]
where each irreducible representation
$\sigma \in \Pi(\overline{K_0})$ occurs at
most with finite multiplicity. For the basis above, one must pick a element
in each one of these ``$\overline{K_0}$-types''.
\end{example}

\begin{example}\label{ex:5.20}
Let $G=SL(2),\;P=B$. Recalling the determination of
${\mathcal{B}}_B(\pi)_\chi $ above, we have 
\[
K_{B,\chi }(x,y)
=\int_{\Pi _u(\overline{A(\aaa )})}
\sum_{\phi \in {\mathcal{B}}_B(\pi)_\chi }
I_B(\pi ,f)\phi (x)d\pi \cdot \delta _B(x)\overline{\phi (y)}
\]
\end{example}

\begin{example}\label{ex:5.21}
Let $G=SL(2),\;P=G$. Recalling the determination
of ${\mathcal{B}}_B(\pi)_\chi $ above, we have 
\begin{eqnarray*}
K_{G,\chi }(x,y)
=\int_{\Pi _u(\overline{G(\aaa )})}
\sum_{\phi \in {\mathcal{B}}_G(\pi)_\chi }
E(x,I_G(\pi ,f)\phi )d\pi \cdot \overline{E(y,\phi )}+ \\
\int_{\Pi _u(\overline{A(\aaa )})}
\sum_{\phi \in {\mathcal{B}}_B(\pi )_\chi }
I_B(\pi,f)\phi (x)d\pi \cdot \delta _B(x)\overline{\phi (y)}.
\end{eqnarray*}
\end{example}

\begin{example}\label{ex:5.22}
Let $G=SL(2)$. We have 
\begin{eqnarray*}
k_\chi ^T(x,f)
=K_{G,\chi }(x,x)-\sum_{\delta \in B(\qqq )\backslash G(\mathbb{Q})}
K_{B,\chi }(\delta x,\delta x)
\tau _B^{\wedge }(H_0(\delta x)-T) \\
=\int_{\Pi _u(\overline{G(\aaa )})}
\sum_{\phi \in {\mathcal{B}}_G(\pi )_\chi}
E_G(x,I_G(\pi ,f)\phi )d\pi \cdot \overline{E_G(x,\phi )}+ \\
+\int_{\Pi _u(\overline{A(\aaa )})}
\sum_{\phi \in {\mathcal{B}}_B(\pi )_\chi }
I_B(\pi,f)\phi (\delta x)d\pi \cdot \delta _B(\delta x)
\overline{\phi (\delta x)}(1-\tau _B^{\wedge }(H(x)-T)) \\
-\sum_{ \delta \in B(\qqq )\backslash G(\qqq ), \delta \neq 1 }
\int_{\Pi _u(\overline{A(\aaa )})}
\sum_{\phi \in {\mathcal{B}}_B(\pi )_\chi }I_B(\pi,f)
\phi (\delta x)d\pi \cdot \delta _B(\delta x)
\overline{\phi (\delta x)}\tau _B^{\wedge }(H(x)-T)
\end{eqnarray*}
\end{example}

Let 
\begin{equation}\label{eq:5.21}
K_\chi ^T(f)=\int_{G(\qqq )\backslash \overline{G(\aaa )}}
\Lambda _2^TK_\chi (x,x)dx.
\end{equation}
Define 
\begin{equation}\label{eq:5.22}
J_\chi ^T(x,y)=\sum_Pn(A)^{-1}
\int_{\Pi _u(\overline{M(\aaa )})}
\sum_{\phi \in {\mathcal{B}}_P(\pi )_\chi }
E(x,I_P(\pi ,f)\phi )\overline{\Lambda ^TE(y,\phi )}d\pi
\end{equation}
and $j_\chi ^T(x,f)=J_\chi ^T(x,x)$. Let 
\[
J_\chi ^T(f)
=\int_{G(\qqq )\backslash \overline{G(\aaa )}}
J_\chi ^T(x,x)dx 
\]
and let
\begin{equation}\label{eq:5.23}
J_\chi (f)=J_\chi ^T(f)|_{T=0}.
\end{equation}

\begin{lemma}\label{lem:5.23}
The integrals $K_\chi ^T(f),J_\chi ^T(f)$ are absolutely convergent and $%
K_\chi ^T(f)=J_\chi ^T(f)$ if $T$ is sufficiently regular.
\end{lemma}

\begin{remark}\label{rem:5.24}
The equality above is the metaplectic analog of \cite[Lemma 2.4]{A3}
for $SL(2)$.
\end{remark}

proof in the case $G=SL(2)$: The integral 
\[
K_\chi ^T(f)
=\int_{G(\qqq )\backslash \overline{G(\aaa )}}
\Lambda _2^TK_\chi (x,x)dx. 
\]
is absolutely convergent since the integrand is compactly
supported on $G(\qqq )\backslash \overline{G(\aaa )}$.

Let $F(P_1,P_2)$ denote the subset of elements of
$P_1(\qqq )\backslash P_2(\qqq )$ which do not belong
to for some parabolic $P$ with $P_1\subset
P\subset P_2,\;P\neq P_2$. Then 
\[
F(P_1,P_2)=\left\{ 
\begin{array}{c}
\{1\},\;\;\;\;P_1=P_2, \\ 
B(\qqq )\backslash G(\qqq )-\{1\},\;\;P_1=B,\;P_2=G,
\end{array}
\right. 
\]
and, from \cite[p. 105]{A3}, 
\begin{eqnarray*}
&&\int_{G(\qqq )\backslash \overline{G(\aaa )}}
k_\chi ^T(x,f)dx \\
&&=\sum_{P_1\subset P_2}(-1)^{\dim (A_2/Z)}
\int_{P_1(\qqq )\backslash \overline{G(\aaa )}}
\sum_{\gamma \in F(P_1,P_2)}\sigma _1^2(H(x)-T)
\Lambda_2^{T,P_1}K_{P_1,\chi }(\gamma x,x)dx \\
&&=\int_{G(\qqq )\backslash \overline{G(\aaa )}}
\sigma _G^G(H(x)-T)\Lambda_2^{T,G}K_{G,\chi }(x,x)dx\\
&&\;\;\;\;\;(the\ P_1=P_2=G\ term) \\
&&+\int_{B(\qqq )\backslash \overline{G(\aaa )}}
\sigma _B^G(H(x)-T)
\sum_{\gamma\in B(\qqq )\backslash G(\qqq )-\{1\}}
\Lambda _2^{T,B}K_{B,\chi }(\gamma x,x)dx\\
&&\;\;\;\;\;(the\ P_1 =B,\;P_2=G\ term) \\
&&-\int_{B(\qqq )\backslash \overline{G(\aaa )}}
\sigma _B^B(H(x)-T)\Lambda_2^{T,B}K_{B,\chi }(x,x)dx
\;\;\;\;\;(P_1 =P_2=B)
\end{eqnarray*}
Each of these integrals converge absolutely since the integral is continuous
and compactly supported. $\Box $

\begin{lemma}\label{lem:5.25}
$\sum_{ \chi =(M,\rho )\in \mathcal{X}, M\neq G }
\int_{G(\qqq )\backslash \overline{G(\aaa )}}
|j_\chi ^T(x,f)|dx<\infty $.
\end{lemma}

proof in the case $G=SL(2)$: Since
$K_P(x,y)=\sum_{\chi \in {\mathcal{X}}}K_{P,\chi }(x,y)$
is absolutely and uniformly convergent on compacta in
$x,y $, $K_P(x,y)$ is continuous in $x,y$. The same
reasoning as in the
previous lemma therefore applies here as well. $\Box $

\begin{lemma}\label{lem:5.26}
For $\chi \in \mathcal{X}$, we have 
\[
\int_{G(\qqq )\backslash \overline{G(\aaa )}}
j_\chi^T(x,f)dx=\sum_Pn(A)^{-1}
\int_{\Pi _u(\overline{M(\aaa )})}
\sum_{\phi \in {\mathcal{B}}_P(\pi )_\chi }
\Lambda ^TE(x,I_P(\pi ,f)\phi )\overline{\Lambda ^TE(y,\phi )}d\pi .
\]
\end{lemma}

proof: This follows from Lemma~\ref{lem:5.23} above on the
equality of $J_\chi ^T(f)$ and $K_\chi ^T(f)$, along with the
metaplectic analog of the fact that
fact $\Lambda^T\circ \Lambda ^T=\Lambda ^T$
\cite[Corollary 1.2]{A3} and that $\Lambda ^T$ is
self-adjoint \cite[Lemma 1.3]{A3} (see Lemma~\ref{lem:5.13} above). $\Box $

Let $L_P\subset {{\mathfrak{a}}}_P^G$ denote the integral
lattice generated by the
co-roots $\{\alpha ^{\vee }\;|\;\alpha \in \Delta _P\}$. For $\phi \in 
{\mathcal{H}}_{P_1}^0(\pi ),
\;s\in \Omega ({{\mathfrak{a}}}_1,{{\mathfrak{a}}}_2),\varsigma
\in {{\mathfrak{a}}}_{1,\mathbb{C}}^{*}$, let 
\[
M(s,\pi ,\varsigma )\phi =(M(s,\pi _\varsigma )\phi _\varsigma
)_{-s\varsigma }. 
\]
\begin{theorem}\label{thm:5.27}
If $\chi \in \mathcal{X}$ is an unramified class then (\cite[section 4]{A3}) 
\[
J_\chi ^T(f)=\sum_{P\in \mathcal{P}_\chi }
n(A)^{-1}\int_{\Pi _u(\overline{M(\aaa )})}
tr(M_P^T(\pi )_\chi I_P(\pi ,f)_\chi )d\pi ,
\]
where 
\[
M_P^T(\pi )_\chi =vol({{\mathfrak{a}}}_P^G/L_P)\lim_{\varsigma \rightarrow
0}\sum_{P_2\in \mathcal{P}_\chi }
\sum_{s\in \Omega ({{\mathfrak{a}}},{{\mathfrak{a}}}_2)}
\frac{e^{(s\varsigma )(T)}M(s,\pi )^{-1}M(s,\pi ,\varsigma )}{\prod_{\alpha
\in \Delta _2}(s\varsigma )(\alpha ^{\vee })}.
\]
\end{theorem}

For a similar but more general result (in the non-metaplectic case),
see \cite[Theorem 8.2]{A14}, which does not assume that $\chi$ is unramified.

\begin{example}\label{ex:5.28} For $G=SL(2)$, each $(G,\pi)\in\chi\in{\mathcal{X}}$ is
unramified. Moreover, $(A,\mu)\in\chi\in{\mathcal{X}}$ is unramified
if and only if $\mu^w\not=\mu$.
($P=G,\chi =(G,\rho )$) $M_G^T(\pi )_\chi =1$

($P=B,\chi =(G,\rho )$) $M_B^T(\pi )_\chi =1$

($P=B,\chi =(A,\rho )$ unramified) 
\[
\begin{array}{c}
M_B^T(\pi )_\chi \\
=vol({\mathbb{R}}/{\mathbb{Z}})
\lim_{\varsigma \rightarrow 0}
\left[
\frac{e^{\varsigma (T)}
M(1,\pi )^{-1}M(1,\pi ,\varsigma )}{\varsigma (\alpha^{\vee })}
+\frac{e^{-\varsigma (T)}
M(s,\pi )^{-1}M(s,\pi ,\varsigma )}{-\varsigma (\alpha^{\vee })}
\right]\\
=M(s,\pi)^{-1}
\lim_{\varsigma \rightarrow 0}
{M(s,\pi)-e^{-2\varsigma (T)}M(s,\pi,\varsigma)
\over \varsigma(\alpha^\vee)} ,
\end{array}
\]
where $\alpha $ denotes the unique positive root.
We identify ${\mathfrak{a}}^*$ with ${\mathbb{R}}$, ${\mathfrak{a}}$ with
${\mathbb{R}}$, and
${\mathfrak{a}}^*\times {\mathfrak{a}}\rightarrow {\mathbb{R}}$ with
the map $(x,y)\longmapsto 2xy$
of real numbers. Thus, at a function $\phi$
on $K$, the above limit becomes 
\[
\begin{array}{c}
M(s,\pi)^{-1}
\lim_{\varsigma \rightarrow 0}
{M(s,\pi)\phi(k)-e^{-2\varsigma (T)}e^{\varsigma (H(k))}
M(s,\pi_\varsigma)\phi_\varsigma(k)
\over \varsigma(\alpha^\vee)} \\
=M(s,\pi)^{-1}
\lim_{\varsigma \rightarrow 0}
{M(s,\pi)\phi(k)-e^{-2\varsigma (T)}M(s,\pi_\varsigma)\phi_\varsigma
\over \varsigma(\alpha^\vee)} \\
=-M(s,\pi)^{-1}M'(s,\pi)\phi(k)+4 \phi(k) T,
\end{array}
\]
where the prime denotes differentiation with respect to
$\varsigma$. Thus
\[
M_B^T(\pi )_\chi =
-M(s,\pi)^{-1}M'(s,\pi)+4\cdot I\cdot T,
\]
where $I$ denotes the identity. (There appears to be a sign
mistake here - the minus sign perhaps differs with some
calculations of Arthur. However, the sign it gives for the
corresponding expression in the trace formula
seems to agree with the corresponding sign in \cite{LL}.)

For
examples and properties of more general limits, see \cite[section 6]{A6}.
\end{example}


\subsection{The trace formula}

The geometric side of the trace formula is the sum of the following terms:

(a) the sum over classes ${\mathcal{O}}\subset G(\qqq )$
such that ${\mathcal{O}}\cap Z(G(\qqq ))\not= \emptyset$ of
\begin{equation}\label{eq:5.24}
J_{\mathcal{O}}(f)=
a^G(S,z)f(z)+{1\over 2}a^A(S,z)J_A(z,f),
\end{equation}
where $J_A$ and $a^M$ are given in \S\S\ref{ssec:JMu}--\ref{ssec:explicit} (these include the
``nilpotent terms''),

(b) the sum over elliptic regular classes ${\mathcal{O}}\subset G(\qqq )$
of
\begin{equation}\label{eq:5.25}
J_{\mathcal{O}}(f)=
a^G(S,\gamma)J_G(\gamma,f),
\end{equation}
where $\gamma\in {\mathcal{O}}$ and
$J_G$ (an orbital integral) and $a^G$ (a volume)
are given in \S\S\ref{ssec:JMu}--\ref{ssec:explicit},

(c) the sum over ``hyperbolic'' regular classes
${\mathcal{O}}\subset G(\qqq )$
of
\begin{equation}\label{eq:5.26}
J_{\mathcal{O}}(f)=
a^A(S,\gamma)J_A(\gamma,f),
\end{equation}
where $J_A$ (a weighted orbital integral) and $a^A$ (a volume)
are given in \S\S\ref{ssec:JMu}--\ref{ssec:explicit}.

The spectral side of the trace formula is the sum of the following terms:

(a) the sum over the cuspidal classes $\chi\in {\mathcal{X}}(G)$ of
\begin{equation}\label{eq:5.27}
J_\chi(f)=
{\rm tr}(\pi(f)),\ \ \ (\pi,G)\in\chi,
\end{equation}

(b) the sum over classes $\chi\in {\mathcal{X}}-{\mathcal{X}}(G)$ of
\begin{equation}\label{eq:5.28}
J_\chi(f)=
{1\over 2}\int_{\Pi_u(\overline{A(\aaa _S)})}
{\rm tr}(M_A(\pi)_\chi I_B(\pi,f)_\chi)d\pi,
\end{equation}
where $M_A(\pi)_\chi =M_A^T(\pi)_\chi|_{T=0}$ is
(essentially a logarithmic derivative of an intertwining operator)
given by Example~\ref{ex:5.28}.


